\documentclass[11pt]{article}
\usepackage[margin=1in]{geometry}
\usepackage{amsmath, amssymb}
\usepackage{booktabs}
\usepackage{graphicx}
\usepackage{xcolor}
\usepackage{hyperref}
\usepackage{microtype}
\usepackage{parskip}
\usepackage{caption}
\usepackage{subcaption}
\usepackage{enumitem}
\usepackage{amsthm}
\newtheorem{definition}{Definition}

\hypersetup{colorlinks=true, linkcolor=blue, citecolor=blue, urlcolor=blue}

\title{\textbf{Detecting Cytokine Cascade Direction\\
       from Latent Geometry and Training Dynamics}\\[0.5em]
       \large Methods, Validation, and Comparison}
\author{Yam Arieli \\ Mornitzan Lab}
\date{May 2026}

\begin{document}
\maketitle
\tableofcontents
\bigskip

%=======================================================================
\section{Biological Motivation}
%=======================================================================

Cytokine cascades are directional: cytokine $A$ activates cells of type $X$, which produce
cytokine $B$ as a downstream effector. The goal is to infer these directed relationships ---
and identify the relay cell type $X$ --- from single-cell RNA-seq data, without prior knowledge
of which cytokines cascade into which.

The key observation is the following. In the Oesinghaus et al.\ dataset, cells are cultured
in pseudo-tubes: closed bags of $\approx\!480$ PBMCs treated with a single cytokine for 24
hours. If cascade $A\!\to\!B$ exists, cells inside an $A$-treated tube will have begun
responding to \emph{endogenously produced} $B$ --- because cells that were directly activated
by $A$ start secreting $B$ into the same tube. As a result, an $A$-tube contains a mixed
signal: a strong direct response to $A$ and a weaker secondary response to $B$.

Crucially, the reverse is not true: a $B$-treated tube contains no endogenously produced $A$.
This \textbf{directional asymmetry} --- $A$-tubes look partially like $B$-tubes, but not
vice versa --- is the cascade signal we aim to detect.

%=======================================================================
\section{System Overview}
%=======================================================================

The full pipeline is illustrated in Figure~\ref{fig:pipeline}. A residual MLP encoder
maps individual cells to 128-dimensional embeddings; an attention module aggregates these
into a bag-level representation; a linear classifier predicts the cytokine label. The
cascade detection methods operate on the encoder embeddings after applying the PBS-RC
transform described in Section~\ref{sec:pbsrc}.

\begin{figure}[h]
  \centering
  \includegraphics[width=\linewidth]{fig1_pipeline.png}
  \caption{%
    \textbf{AB-MIL Cytokine Cascade Inference Pipeline.}
    Each pseudo-tube ($N \times 4000$ genes) is passed through an InstanceEncoder
    (MLP, $h_i \in \mathbb{R}^{128}$), attention aggregation, and a 91-class
    classifier. The PBS-RC transform (red box) subtracts the cell-type-specific
    resting-state centroid from each embedding before computing cascade scores.
    The green arrow shows where asymmetry scores are read off from the
    attention-weighted centroid trajectories.}
  \label{fig:pipeline}
\end{figure}

%=======================================================================
\section{Mathematical Setting}
%=======================================================================

\textbf{Cells and expression.} Let $\mathcal{I}$ be the set of all cells. Each cell
$i \in \mathcal{I}$ has a gene expression vector $x_i \in \mathbb{R}^G$ (log-normalised,
$G=4{,}000$ HVGs), a cytokine label $c(i) \in \mathcal{C}$ (where $C_0$ denotes PBS), a
cell-type label $\tau(i) \in \mathcal{T}$ ($M=18$ annotated PBMC types), and a donor label
$d(i) \in \mathcal{D}$ ($|\mathcal{D}|=12$).

\textbf{Encoder.} A residual MLP $f_\theta:\mathbb{R}^G\!\to\mathbb{R}^d$ ($d=128$) maps
each cell to a dense embedding $h_i = f_\theta(x_i)$. $f_\theta$ is pre-trained by
cell-type classification (Stage~1); it has never seen cytokine labels, so all cytokine
signal in $h_i$ is emergent.

\textbf{Donors.} The dataset is split at the donor level: $\mathcal{D} =
\mathcal{D}_\mathrm{train} \cup \mathcal{D}_\mathrm{val}$, with
$\mathcal{D}_\mathrm{train} = \{\text{D1, D4--D12}\}$ and
$\mathcal{D}_\mathrm{val} = \{\text{D2, D3}\}$. All centroids are computed on
$\mathcal{D}_\mathrm{train}$ only.

\textbf{Cell sets.} For cytokine $A\in\mathcal{C}$ and cell type $T\in\mathcal{T}$:
$\mathcal{S}_{A,T} = \{i\in\mathcal{I}: c(i)=A,\,\tau(i)=T,\,d(i)\in\mathcal{D}_\mathrm{train}\}$.

%=======================================================================
\section{The PBS-RC Transform}
\label{sec:pbsrc}
%=======================================================================

\subsection{The Composition Problem}

The dominant axis of variation in $\mathbb{R}^d$ is cell type: every NK cell lands near
the NK cluster regardless of cytokine treatment. Figure~\ref{fig:realistic_embedding}
illustrates this: in raw embedding space, cytokine centroids nearly coincide and
cell-type identity overwhelms any cytokine-induced displacement. After the PBS-RC
transform, all cell types are centred at their own resting state and within-cluster
displacements become visible.

\begin{figure}[h]
  \centering
  \includegraphics[width=\linewidth]{fig_realistic_embedding.png}
  \caption{%
    \textbf{PBS-RC resolves the composition confound.}
    \textbf{(a)} Raw embedding space: cell-type identity dominates; NK, Mono, and T
    cells from IL-12, IFN-$\gamma$, and PBS tubes all overlap within their respective
    clusters. IL-12 and IFN-$\gamma$ tube centroids are indistinguishable at this scale.
    \textbf{(b)} PBS-RC space $\tilde{h}_i = h_i - \boldsymbol{\mu}_{\mathrm{PBS},\tau(i)}$:
    all cell types are centred at the origin; NK cells in IL-12 tubes are visibly displaced
    toward IFN-$\gamma$ space (cascade signal). The bias measurement along $\hat{u}_{A\to B}$
    is shown.}
  \label{fig:realistic_embedding}
\end{figure}

A naïve comparison of cytokine centroids $\bar{h}_A$ conflates cell-state changes with
cell-type composition differences between conditions. Even the within-tube displacement
$(\boldsymbol{\mu}_{A,T} - \boldsymbol{\mu}_A)$ fails without PBS subtraction: if
cytokine $A$ causes NK proliferation, $\boldsymbol{\mu}_A$ shifts toward the NK cluster
even if no NK cell changed its transcriptional state, producing a spurious bias. The
PBS-RC transform resolves this.

\subsection{PBS-RC Embedding}

\begin{definition}[PBS reference centroid]
For each cell type $T\in\mathcal{T}$, the PBS reference centroid is
\[
  \boldsymbol{\mu}_{\mathrm{PBS},T}
  = \frac{1}{|\mathcal{S}_{C_0,T}|}\sum_{i\in\mathcal{S}_{C_0,T}} h_i \in \mathbb{R}^d.
\]
\end{definition}

\begin{definition}[PBS-RC embedding]
The PBS-Relative Centroid embedding of cell $i$ is
$\tilde{h}_i = h_i - \boldsymbol{\mu}_{\mathrm{PBS},\tau(i)} \in \mathbb{R}^d$.
\end{definition}

By construction, $\tilde{h}_i = \mathbf{0}$ whenever cell $i$ occupies the same position
as the average resting-state cell of type $\tau(i)$. Because the subtraction is
cell-type-specific, it genuinely reshapes the geometry rather than shifting all cells
by the same global offset.

\subsection{Cytokine Centroids in PBS-RC Space}

For cytokine $A\in\mathcal{C}\setminus\{C_0\}$:
\[
  \boldsymbol{\mu}_A = \frac{1}{|\mathcal{S}_A|}\sum_{i\in\mathcal{S}_A}\tilde{h}_i,
  \qquad
  \boldsymbol{\mu}_{A,T} = \frac{1}{|\mathcal{S}_{A,T}|}\sum_{i\in\mathcal{S}_{A,T}}\tilde{h}_i.
\]
The deviation $(\boldsymbol{\mu}_{A,T} - \boldsymbol{\mu}_A)$ measures how much cell type
$T$ is displaced within $A$-tubes relative to the average $A$-tube cell.
Figure~\ref{fig:concept} summarises the geometric interpretation.

\begin{figure}[h]
  \centering
  \includegraphics[width=\linewidth]{fig_geometry_concept.png}
  \caption{%
    \textbf{Cascade detection via cell-state geometry.}
    \textit{Left:} Raw embedding space. Cell-type identity dominates; cytokine tube
    centroids nearly coincide.
    \textit{Right:} PBS-RC space. All cell types are centred at the origin; the
    observed IL-12 centroid $c_A$ is displaced toward the IFN-$\gamma$ centroid $c_B$,
    whereas the non-cascading direction shows no displacement. The asymmetry score
    $\mathrm{ASYM}(A\!\to\!B) = \mathrm{bias}(A,B) - \mathrm{bias}(B,A) > 0$
    is evidence for cascade $A \to B$.}
  \label{fig:concept}
\end{figure}

Figure~\ref{fig:cascade_embedding} shows the same transform applied to real encoder
embeddings.

\begin{figure}[h]
  \centering
  \includegraphics[width=\linewidth]{fig_cascade_embedding.png}
  \caption{%
    \textbf{PBS-RC reveals within-cluster cytokine displacements.}
    \textit{Left:} Raw encoder space. Cell-type identity dominates; IL-12 and IFN-$\gamma$
    tube centroids nearly coincide, and the NK displacement $\boldsymbol{\mu}_{A,\mathrm{NK}}
    - \boldsymbol{\mu}_A$ spans the whole space.
    \textit{Right:} After the PBS-RC transform, cytokine centroids separate, and NK cells
    in IL-12 tubes are visibly displaced toward IFN-$\gamma$ space (bias~$= +0.09$),
    while Monocytes show no displacement (bias~$= -0.05$).}
  \label{fig:cascade_embedding}
\end{figure}

%=======================================================================
\section{Static Cascade Direction Methods}
%=======================================================================

\subsection{Directional Bias}

For cytokines $A, B \in \mathcal{C}\setminus\{C_0\}$ with $A\neq B$, let
\[
  \mathbf{u}_{A\to B} = \frac{\boldsymbol{\mu}_B - \boldsymbol{\mu}_A}
                             {\|\boldsymbol{\mu}_B - \boldsymbol{\mu}_A\|_2}
\]
be the unit vector pointing from $A$'s centroid to $B$'s centroid. The
\textbf{directional bias} of cell type $T$ in $A$-tubes toward $B$ is:
\[
  \mathrm{bias}(A\!\to\!B, T)
  = \bigl(\boldsymbol{\mu}_{A,T} - \boldsymbol{\mu}_A\bigr)\cdot\mathbf{u}_{A\to B}
  \in \mathbb{R}.
\]

Figure~\ref{fig:bias_explained} shows a worked 2D example: $\mathrm{bias}(\mathrm{IL\text{-}12}
\!\to\!\mathrm{IFN\text{-}\gamma}) = 0.80$ and $\mathrm{bias}(\mathrm{IFN\text{-}\gamma}
\!\to\!\mathrm{IL\text{-}12}) = 0.625$, giving $\mathrm{ASYM} = 0.175 > 0$.

\begin{figure}[h]
  \centering
  \includegraphics[width=0.65\linewidth]{fig_bias_explained.png}
  \caption{%
    \textbf{Directional bias is an asymmetric projection.}
    IL-12 centroid $(1.0,\,0.8)$ and IFN-$\gamma$ centroid $(0.0,\,1.0)$ in a
    2D PBS-RC space. The blue vector is IL-12's deviation projected onto $\hat{u}_{A\to B}$
    (bias $= 0.80$); the red vector is IFN-$\gamma$'s projection onto the same axis
    (bias $= 0.625$). The asymmetry $0.175 > 0$ indicates IL-12 $\to$ IFN-$\gamma$.
    The key property: bias is a projection, not a distance --- it is not symmetric.}
  \label{fig:bias_explained}
\end{figure}

\subsection{Why This Captures Cascade Direction}

Suppose cascade $A \to B$ exists via relay cell type $X$:
\begin{enumerate}[noitemsep]
  \item Cytokine $A$ is added to the pseudo-tube.
  \item Cells of type $X$ are primary responders: directly activated by $A$, they begin
        secreting $B$ autocrinally. Their state reflects \emph{both} $A$ and $B$.
  \item All other cells receive $B$ paracrinally and shift toward $B$-space.
\end{enumerate}
$X$ receives a double signal and shifts \emph{more} than the average tube cell:
$\mathrm{bias}(A\!\to\!B, X) > 0$.
In a $B$-treated tube there is no endogenous $A$; $X$ responds to $B$ like every other
cell type, so $\mathrm{bias}(B\!\to\!A, X) \approx 0$.

\subsection{Legacy Asymmetry Score and Its Structural Problems}

The legacy aggregate score was:
\[
  \mathrm{ASYM}(A\!\to\!B) = \max_{T\in\mathcal{T}}\bigl[
    \mathrm{bias}(A\!\to\!B,T) - \mathrm{bias}(B\!\to\!A,T)
  \bigr].
\]
The $\max$ over cell types is motivated by the fact that a cascade acts through a specific
relay; the argmax directly identifies $X^*$.

Despite being well-motivated, the subtraction introduces two structural defects.

\paragraph{Contamination.}
Since $\hat{u}_{B\to A} = -\hat{u}_{A\to B}$, expanding the subtraction gives:
\[
  \mathrm{bias}(A\!\to\!B, T) - \mathrm{bias}(B\!\to\!A, T)
  =
  \underbrace{(\boldsymbol{\mu}_{A,T} - \boldsymbol{\mu}_A)\cdot\hat{u}_{A\to B}}_{\text{forward (wanted)}}
  +
  \underbrace{(\boldsymbol{\mu}_{B,T} - \boldsymbol{\mu}_B)\cdot\hat{u}_{A\to B}}_{\text{contaminant}}.
\]
A cell type $T$ that responds strongly \emph{directly} to $B$ inflates
$\mathrm{ASYM}(A\!\to\!B)$ even when no cascade through $T$ exists.

\paragraph{Forced antisymmetry.}
$\mathrm{ASYM}(A\!\to\!B) = -\mathrm{ASYM}(B\!\to\!A)$ identically. The method cannot
distinguish a genuine directed cascade from an algebraic sign flip.

Figure~\ref{fig:pbs_rc_diagram} illustrates the method and its failure modes.

\begin{figure}[h]
  \centering
  \includegraphics[width=\linewidth]{fig_pbs_rc_diagram.png}
  \caption{%
    \textbf{PBS-RC geometric method for cascade detection (four-panel summary).}
    \textbf{(A)} Biological mechanism: cytokine $A$ directly activates relay cell type $X$
    (NK cells), which secrete $B$ autocrinally.
    \textbf{(B)} PBS-RC space: NK cells in $A$-tubes are strongly displaced toward
    $B$ (bias~$= 0.70$); Monocytes barely deviate (bias~$= 0.05$).
    \textbf{(C)} Asymmetry test: large forward bias minus small reverse bias gives
    $\mathrm{ASYM}(A\!\to\!B) = 0.65$, identifying NK as the relay $X^*$.
    \textbf{(D)} Why $\boldsymbol{\mu}_{A,T} - \boldsymbol{\mu}_{B,T}$ is the wrong
    formula: when a cascade exists, NK in $A$-tubes resembles NK in $B$-tubes
    (small $|\delta|$, invisible), while Monocytes differ for unrelated reasons
    (large $|\delta|$, spurious). The naïve formula inverts the ranking.}
  \label{fig:pbs_rc_diagram}
\end{figure}

Figure~\ref{fig:asym_heatmap} shows the full 91$\times$91 legacy asymmetry matrix.

\begin{figure}[h]
  \centering
  \includegraphics[width=\linewidth]{fig_asym_heatmap.png}
  \caption{%
    \textbf{Full legacy asymmetry matrix (91$\times$91 ordered pairs).}
    $\mathrm{ASYM}(A\!\to\!B) = \mathrm{bias}(A,B) - \mathrm{bias}(B,A)$. Red~$=$
    $A$ likely upstream of $B$; Blue~$=$ $B$ likely upstream. Seed~42, late 30\%
    of training. Several cytokines (e.g.\ IL-27, IFN-$\beta$) show consistent
    asymmetry rows, suggesting they are common cascade sources. The anti-symmetric
    structure (each row is the sign-flip of the corresponding column) is a structural
    artefact of the legacy formula.}
  \label{fig:asym_heatmap}
\end{figure}

\subsection{Refined Static Pipeline: Per-Donor Wilcoxon}
\label{sec:refined}

The refined pipeline keeps the same geometric intuition but replaces the contaminated
subtraction with two fully independent one-sided statistical tests.

\paragraph{Step 1 --- Per-donor PBS-RC.}
For every cell type $T$ and training donor $d$:
$\tilde{h}_i = h_i - \boldsymbol{\mu}_{\mathrm{PBS},\tau(i)}^{(d_i)}$,
using donor-specific PBS centroids.

\paragraph{Step 2 --- Per-donor forward projection.}
\[
  b_\mathrm{fwd}^{(d)}(A\!\to\!B, T)
  = \bigl(\tilde{\boldsymbol{\mu}}_{A,T}^{(d)} - \tilde{\boldsymbol{\mu}}_A\bigr)
    \cdot \hat{u}_{A\to B},
\]
where $\tilde{\boldsymbol{\mu}}_A = \frac{1}{|\mathcal{D}_\mathrm{train}|}\sum_d
\tilde{\boldsymbol{\mu}}_A^{(d)}$.
No reverse term is subtracted anywhere. The ``reverse'' of $(A,B)$ is the independent
forward test for $(B,A)$.

\paragraph{Step 3 --- Statistical testing.}
One-sided Wilcoxon signed-rank test across $d\in\mathcal{D}_\mathrm{train}$ (null:
median $\le 0$) for each $(A,B,T)$. Two-stage correction: Bonferroni across cell types
within each ordered pair (relay search), then Benjamini--Hochberg FDR ($\alpha = 0.05$)
across all $K(K-1) = 8{,}010$ ordered pairs.

\paragraph{Cascade call.}
$\text{fwd\_sig} = \exists\,T: p_\mathrm{Bonf}(A\!\to\!B,T)\le\alpha$,
$\text{rev\_sig} = \exists\,T: p_\mathrm{Bonf}(B\!\to\!A,T)\le\alpha$,
producing one of: \texttt{A->B}, \texttt{B->A}, \texttt{shared}, or \texttt{none}.

%=======================================================================
\section{Centroid Trajectory Method}
\label{sec:trajectory}
%=======================================================================

The static methods use a single frozen-encoder snapshot. An alternative approach
exploits the \emph{dynamics} of fine-tuning: if cascade $A\!\to\!B$ exists, the
cell-type centroids in $A$-tubes should systematically drift toward $B$'s centroid
as the model is fine-tuned to discriminate cytokine labels.

\subsection{Training Protocol}

Figure~\ref{fig:ct_method} shows the three-stage training protocol and the trajectory
concept.

\begin{figure}[h]
  \centering
  \includegraphics[width=\linewidth]{fig_ct_method.png}
  \caption{%
    \textbf{Training protocol and centroid trajectory concept.}
    \textit{Left:} Three stages. Stage~1: encoder pre-training on cell-type labels
    (50 epochs, lr~$= 0.01$). Stage~2a: frozen-encoder MIL warmup (20 epochs,
    lr~$= 0.001$). Stage~2b/3: full fine-tuning of both encoder and MIL weights
    (100 epochs, lr~$= 0.001$) with attention-weighted centroid snapshots logged every
    10 epochs. Small LR chosen for smooth trajectories, not spiky jumps.
    \textit{Right:} Trajectory concept. If $A\!\to\!B$ cascade exists, attention-weighted
    centroids of relay cell type $T$ in $A$-tubes (green dots, $t_1\to t_2\to t_3$) drift
    toward $\boldsymbol{\mu}_B$ along the $\hat{u}_{A\to B}$ axis, giving a positive
    $b_\mathrm{fwd}$ slope.}
  \label{fig:ct_method}
\end{figure}

At each logged epoch $t$, the \textbf{attention-weighted centroid} for cytokine $C$,
cell type $T$, donor $d$ is:
\[
  \tilde{\boldsymbol{\mu}}_{C,T}^{(d)}(t)
  = \frac{\sum_{i \in \mathcal{S}_{C,T}^{(d)}} a_i(t)\,\tilde{h}_i(t)}
         {\sum_{i \in \mathcal{S}_{C,T}^{(d)}} a_i(t)},
\]
where $a_i(t)$ are the MIL attention weights and $\tilde{h}_i(t)$ is the PBS-RC embedding
at epoch $t$. The trajectory score for pair $(A,B)$, cell type $T$, donor $d$ is the
OLS slope of $\bigl\{b_\mathrm{fwd}^{(d)}(A\!\to\!B, T, t)\bigr\}_{t=10,20,\ldots,100}$.
A positive slope means relay cell type $T$ drifts progressively toward $B$'s centroid.

\subsection{Generalization Check}

Figure~\ref{fig:ct_accuracy} confirms the model learns cytokine-specific programs that
generalise to held-out donors.

\begin{figure}[h]
  \centering
  \includegraphics[width=\linewidth]{fig_ct_accuracy.png}
  \caption{%
    \textbf{Final-epoch classification accuracy across 10 seeds.}
    Train donors (D1, D4--D12): mean P(correct) ranges from $0.57$ to $0.86$.
    Val donors (D2, D3), observer-only (zero gradient updates): P(correct) ranges
    from $0.047$ to $0.079$ --- $4$--$7\times$ above chance ($1/91 \approx 0.011$,
    dotted line). Despite the train/val gap, val accuracy is significantly above
    chance, confirming that the encoder has learnt cytokine-specific programs that
    generalise beyond the training donors.}
  \label{fig:ct_accuracy}
\end{figure}

\subsection{NK Slope Analysis for the Positive Control (IL-12 \texorpdfstring{$\to$}{->} IFN-\texorpdfstring{$\gamma$}{gamma})}

The IL-12 $\to$ IFN-$\gamma$ cascade (mediated by NK cells) is the best-documented
PBMC cascade in the dataset. Figure~\ref{fig:ct_nk_slopes} shows the per-donor NK
trajectory slopes across 10 seeds.

\begin{figure}[h]
  \centering
  \includegraphics[width=\linewidth]{fig_ct_nk_slopes.png}
  \caption{%
    \textbf{NK cell bias-trajectory slope for IL-12 $\to$ IFN-$\gamma$, per donor per seed.}
    Each cell is the OLS slope of $b_\mathrm{fwd}(\mathrm{IL\text{-}12}\!\to\!\mathrm{IFN\text{-}\gamma},
    \mathrm{NK}, d, t)$ over the 10 logged Stage-3 epochs. Positive (blue) $=$ NK cells
    in IL-12 tubes drift progressively toward IFN-$\gamma$ centroid during fine-tuning.
    Seeds split into a \textbf{positive group} (6 seeds, top rows; dashed boundary) and
    a \textbf{negative group} (4 seeds, bottom rows). The positive group shows consistent
    cascade-direction drift across all 10 training donors.}
  \label{fig:ct_nk_slopes}
\end{figure}

Figure~\ref{fig:ct_nk_summary} aggregates the per-donor slopes to the seed level.

\begin{figure}[h]
  \centering
  \includegraphics[width=\linewidth]{fig_ct_nk_summary.png}
  \caption{%
    \textbf{Per-donor NK slopes aggregated by seed (IL-12 $\to$ IFN-$\gamma$).}
    Each dot is one training donor; the horizontal bar is the seed mean. 6/10 seeds
    (blue, left cluster) show consistently positive NK slopes; 4/10 seeds (orange,
    right cluster) show negative slopes. The bimodal split suggests two attractor
    basins in the fine-tuning landscape; the source of this split is under investigation.
    Pair-level benchmark calls aggregate over all available seeds.}
  \label{fig:ct_nk_summary}
\end{figure}

\subsection{PCA Trajectory Visualization}

Figures~\ref{fig:centroid_pca_nk} and~\ref{fig:centroid_pca_all} show the centroid
trajectories in the cascade-relevant 2D PCA space.

\begin{figure}[h]
  \centering
  \includegraphics[width=0.75\linewidth]{centroid_pca_nk_only.png}
  \caption{%
    \textbf{NK centroid trajectory in PCA space (IL-12 $\to$ IFN-$\gamma$, seed~42).}
    Each dot is one logged epoch (light $=$ early, dark $=$ late); open circle $=$
    epoch~10; star $=$ epoch~100. PCA is fitted jointly on both cytokines across all
    cell types. The IL-12 NK centroid (red/orange) and the IFN-$\gamma$ NK centroid
    (blue) move over fine-tuning and converge directionally: the IL-12 centroid drifts
    toward the IFN-$\gamma$ cluster, consistent with the cascade signal.}
  \label{fig:centroid_pca_nk}
\end{figure}

\begin{figure}[h]
  \centering
  \includegraphics[width=\linewidth]{centroid_pca_traj_IL12_IFNgamma.png}
  \caption{%
    \textbf{All-cell-type centroid trajectories in PCA space (IL-12 $\to$ IFN-$\gamma$, seed~42).}
    Each panel shows one of the 18 PBMC cell types. IL-12 centroids (red/orange) and
    IFN-$\gamma$ centroids (blue) show distinct trajectory patterns by cell type: NK,
    NKT, and ILC cell types exhibit strong directional drift consistent with the cascade
    relay; B-cell and Mono populations show more stable or divergent trajectories,
    reflecting the cell-type specificity of the relay mechanism.}
  \label{fig:centroid_pca_all}
\end{figure}

\subsection{Benchmark Recovery}
\label{sec:ct_benchmark}

The trajectory method was evaluated against 11 literature-curated cascade pairs
(9 positive + 2 controls). Figure~\ref{fig:ct_benchmark} shows the result.

\begin{figure}[h]
  \centering
  \includegraphics[width=\linewidth]{fig_ct_benchmark.png}
  \caption{%
    \textbf{Benchmark pair recovery using the centroid trajectory slope signal.}
    All 11 literature-curated cascade pairs $\times$ all 10 seeds are correctly recalled
    (call $\in \{\texttt{A\texttt{->}B},\,\texttt{shared}\}$; BH-FDR $q = 0.018$,
    Bonferroni across cell types). Every cell shows a green checkmark: 100\% recall at
    the pre-specified statistical threshold.}
  \label{fig:ct_benchmark}
\end{figure}

The trajectory method achieves \textbf{11/11 pairs $\times$ 10/10 seeds} (100\% recall)
on the benchmark set. This is the strongest validation result to date, substantially
improving on the static methods (see Section~\ref{sec:benchmark_comparison}).

%=======================================================================
\section{Experimental Setup}
%=======================================================================

\subsection{Batch Construction}

The full dataset contains 12 donors. For each cytokine $c_i$ ($i=0,\ldots,90$),
one tube is selected from donor $(i+\delta)\!\!\mod\!12$. Two mutually exclusive batches:
\begin{itemize}[noitemsep]
  \item \textbf{Batch~0} ($\delta=0$, seeds 1--10, 42, 123).
  \item \textbf{Batch~1} ($\delta=1$, seeds 11--22).
\end{itemize}
Every cytokine is represented by a \emph{different} donor across the two batches, so the
training cells are mutually exclusive.

\subsection{Pre-Registered Controls}

Before inspecting any results, four cytokine pairs were designated as controls
(Table~\ref{tab:controls}). Figure~\ref{fig:controls} shows the PBS-RC rank percentile
for these pairs before and after the transform.

\begin{table}[h]
\centering
\caption{Pre-registered pairs of interest.}
\label{tab:controls}
\small
\begin{tabular}{llll}
\toprule
Direction & Pair & Role & Expected signal \\
\midrule
$A\!\to\!B$ & IL-12 $\to$ IFN-$\gamma$ & Positive control & High score; $A\!\to\!B > B\!\to\!A$ \\
$B\!\to\!A$ & IFN-$\gamma$ $\to$ IL-12  & Positive control (reverse) & Lower score than forward \\
$A\!\to\!B$ & IL-6 $\to$ IL-10          & Negative control & Mid-range, symmetric with reverse \\
$B\!\to\!A$ & IL-10 $\to$ IL-6          & Negative control (reverse) & Similar rank to forward \\
\bottomrule
\end{tabular}
\end{table}

\begin{figure}[h]
  \centering
  \includegraphics[width=\linewidth]{fig2_controls.png}
  \caption{%
    \textbf{Pre-registered control pair ranks (PBS-RC legacy method).}
    \textit{Left:} Rank percentile before (red) and after (green) PBS-RC transform.
    IL-12 $\to$ IFN-$\gamma$ improves from the 60th to the 18th percentile (strong
    signal zone, below the 10th-percentile dotted threshold). IL-6 $\to$ IL-10 moves
    from the 11th to the 49th percentile --- correctly reclassified as a shared-pathway
    pair after PBS-RC removes the composition confound.
    \textit{Right:} Per-seed ranks across the 12 stable seeds. IL-12 $\to$ IFN-$\gamma$
    (blue) is consistently in the upper decile; IL-6 $\to$ IL-10 (orange) fluctuates
    near the 50th percentile.}
  \label{fig:controls}
\end{figure}

%=======================================================================
\section{Results: Static Methods}
%=======================================================================

\subsection{Within-Batch Seed Stability}

A seed is \emph{stable} if its Spearman $\rho$ with the batch ensemble mean is $\ge 0.70$.
Figures~\ref{fig:rho_offset0} and~\ref{fig:rho_offset1} show per-seed $\rho$ for
Batch~0 and Batch~1.

\begin{figure}[h]
  \centering
  \begin{subfigure}[b]{0.48\linewidth}
    \includegraphics[width=\linewidth]{geo_results/rho_plot_offset0.png}
    \caption{Batch~0 ($\delta=0$, 12 seeds).}
    \label{fig:rho_offset0}
  \end{subfigure}
  \hfill
  \begin{subfigure}[b]{0.48\linewidth}
    \includegraphics[width=\linewidth]{geo_results/rho_plot_offset1.png}
    \caption{Batch~1 ($\delta=1$, 12 seeds).}
    \label{fig:rho_offset1}
  \end{subfigure}
  \caption{%
    \textbf{PBS-RC ensemble seed stability (legacy method).}
    Green bars: stable ($\rho \ge 0.70$, dashed line). Red bars: unstable. Batch~0:
    10/12 stable (seeds~3 and~4 unstable). Batch~1: 11/12 stable. Within-batch
    consistency is high for both independent donor assignments.}
  \label{fig:stability_both}
\end{figure}

\subsection{Cross-Batch Reproducibility}

Figure~\ref{fig:cross_offset} shows the cross-batch scatter for the legacy method.
The two batches use entirely disjoint training cells.

\begin{figure}[h]
  \centering
  \includegraphics[width=0.55\linewidth]{geo_results/cross_offset_plot.png}
  \caption{%
    \textbf{Cross-batch asymmetry score agreement (legacy method).}
    Each point is one of the 8,010 ordered pairs. Spearman $\rho = 0.766$
    ($p < 10^{-300}$) between Batch-0 and Batch-1 stable ensembles, computed on
    disjoint training cells. The slight offset from the $y=x$ line reflects a
    donor-offset effect, not a signal difference.}
  \label{fig:cross_offset}
\end{figure}

\paragraph{Top discovered pairs.}
Two cytokines dominate the top of both batch rankings: \textbf{IL-27} (documented
downstream targets include IL-22 and TNF-$\alpha$) and \textbf{IFN-$\beta$} (type-I
interferon with downstream ISG programme activation). Figure~\ref{fig:stability_pairs}
shows the top pairs from Batch~0 and per-seed stability.

\begin{figure}[h]
  \centering
  \includegraphics[width=\linewidth]{fig3_stability_pairs.png}
  \caption{%
    \textbf{Legacy method --- seed stability and top cascade pairs (Batch~0).}
    \textit{Left:} Per-seed Spearman $\rho$ versus the batch ensemble mean. 10/12 seeds
    stable; seeds~3 and~4 unstable (red).
    \textit{Right:} Top pairs with $\ge 6/10$ stable seeds, ranked by mean legacy
    asymmetry score. IL-27 (orange) and IFN-$\beta$ (blue) dominate. Their consistent
    appearance in both disjoint batches is evidence of a biological rather than
    data-artefact signal.}
  \label{fig:stability_pairs}
\end{figure}

\subsection{Legacy vs.\ Refined: Stability and Cross-Batch Comparison}

Both methods were evaluated on the full 24-run corpus (Batch~0 and Batch~1, 12 seeds each).
Figure~\ref{fig:stability_comparison} compares per-seed $\rho$ side by side.

\begin{figure}[h]
  \centering
  \includegraphics[width=\linewidth]{stability_plot.png}
  \caption{%
    \textbf{Within-batch seed stability: Legacy vs.\ Refined.}
    Dashed line: $\rho=0.70$ threshold. Legacy (blue) has five seeds below threshold
    in Batch~0 (seeds 3, 4, 5, 8, 123); Refined (orange) has one (seed~123 only).
    Batch~1 is comparable for both methods (11/12 stable each). The refined method
    improves within-batch stability from 7/12 to 11/12 in Batch~0.}
  \label{fig:stability_comparison}
\end{figure}

Figure~\ref{fig:cross_batch_scatter} shows the cross-batch scatter for both methods
under different pair-set filters.

\begin{figure}[h]
  \centering
  \includegraphics[width=\linewidth]{scatter_plot.png}
  \caption{%
    \textbf{Cross-batch scatter (Batch~0 score, $x$-axis; Batch~1 score, $y$-axis).}
    \textit{Top row:} all 8,010 pairs. Legacy $\rho=0.772$ (left); Refined $\rho=0.815$
    (right). \textit{Bottom row:} top 50\% by each method's own magnitude.
    Legacy (left) loses coherence ($\rho$ drops to $0.62$), indicating contamination
    drives the largest scores. Refined (right) gains it ($\rho$ rises to $0.84$),
    indicating the refined scores carry genuine signal that concentrates at the top.}
  \label{fig:cross_batch_scatter}
\end{figure}

Table~\ref{tab:cross_batch} reports cross-batch $\rho$ and top-50 overlap under four
pair-set filters. Row~3 is the critical diagnostic: restricting to the top 50\% by
legacy magnitude, legacy $\rho$ \emph{drops} to $0.620$ while refined $\rho$ on the
same pairs \emph{rises} to $0.867$. The legacy method's highest-magnitude pairs are its
\emph{least} reproducible.

\begin{table}[h]
\centering
\caption{Cross-batch Spearman $\rho$ and top-50 overlap under four pair-set filters.
         ``Top-50 overlap'' counts how many of the top-50 pairs by Batch-0 score also
         appear in the top-50 of Batch-1. Row~3 is the critical diagnostic.}
\label{tab:cross_batch}
\small
\begin{tabular}{lrrrr}
\toprule
\textbf{Filter} & \textbf{Legacy $\rho$} & \textbf{Refined $\rho$}
               & \textbf{Legacy overlap} & \textbf{Refined overlap} \\
\midrule
All pairs ($n=8{,}010$)                       & 0.772 & 0.815 & 25/50 &  3/50 \\
Top 50\% by refined score ($n=4{,}005$)       & 0.762 & 0.835 & 28/50 & 10/50 \\
Top 50\% by $|$legacy$|$ score ($n=4{,}005$) & 0.620 & 0.867 & 25/50 &  9/50 \\
Active pairs (refined $>0$, $n=7{,}100$)      & 0.772 & 0.813 & 27/50 &  3/50 \\
\bottomrule
\end{tabular}
\end{table}

\subsection{Benchmark Comparison: Static Methods}
\label{sec:benchmark_comparison}

Figure~\ref{fig:benchmark_comparison} evaluates both static methods against the 11
literature-curated benchmark pairs.

\begin{figure}[h]
  \centering
  \includegraphics[width=\linewidth]{benchmark_figure.png}
  \caption{%
    \textbf{Benchmark pair recovery: Legacy vs.\ Refined (static methods).}
    \textbf{(A)} Legacy score: forward vs.\ reverse asymmetry rank percentile for each
    of the 11 benchmark pairs (lower $=$ stronger signal). Colors: correct $A\!\to\!B$
    call (green), shared (orange), wrong (red), negative control (grey), reverse
    direction (blue). IL-12 $\to$ IFN-$\gamma$ is called ``reverse direction'' by the
    legacy score --- a contamination artefact.
    \textbf{(B)} Refined score and dominant cascade call across 24 seeds. Calls shown
    as ``call (count/24)''. Several benchmark pairs show ``none (0/24)'', indicating
    the refined method is conservative at the BH-FDR $\alpha=0.05$ threshold. The
    centroid trajectory method (Figure~\ref{fig:ct_benchmark}) achieves 100\% recall
    on the same pairs.}
  \label{fig:benchmark_comparison}
\end{figure}

%=======================================================================
\section{Limitations and Caveats}
%=======================================================================

\begin{enumerate}[noitemsep]
  \item \textbf{Encoder is cell-type supervised, not cytokine supervised.} Cytokine
        signals are encoded incidentally as deviations within cell-type clusters. If
        cytokine $A$ induces a programme identical to resting-state variation in cell
        type $T$, it will not appear as a displacement in PBS-RC space.

  \item \textbf{Small effective donor sample.} All centroids are computed from
        $|\mathcal{D}_\mathrm{train}|=10$ donors. Pseudo-tubes from the same donor are
        correlated; effective $n=10$, not the number of cells. Statistical power is limited
        and the donor-level Wilcoxon test is the correct unit of analysis.

  \item \textbf{24-hour snapshot.} All tubes are measured at a single 24-hour time point.
        Cytokines with slow induction kinetics may not have produced sufficient secondary
        $B$ for the effect to appear. A positive score is consistent with a cascade but
        could also reflect simultaneous shared programme activation.

  \item \textbf{Bimodal seed split in the trajectory method.} 4/10 seeds show negative
        NK slopes for the IL-12 $\to$ IFN-$\gamma$ positive control
        (Figure~\ref{fig:ct_nk_summary}), suggesting multiple attractor basins in the
        fine-tuning landscape. The source of this split is not yet understood.

  \item \textbf{Top-50 overlap remains low for static methods.} Even under the refined
        method, top-50 cross-batch overlap reaches only 10/50. Monotone improvement
        with filtering suggests a threshold problem rather than absent signal, but
        stricter FDR thresholds or more seeds are needed to verify.
\end{enumerate}

%=======================================================================
\section{Conclusions}
%=======================================================================

\begin{enumerate}[noitemsep]
  \item \textbf{The PBS-RC transform is necessary.} Without it, cell-type composition
        confounds cytokine centroid comparisons. After the transform, within-cluster
        cytokine-induced displacements become visible and interpretable.

  \item \textbf{Static methods are reproducible above threshold.} Within-batch stability
        is high (10--11/12 seeds) and cross-batch $\rho$ exceeds the pre-specified
        $0.70$ threshold on disjoint training data (legacy $0.766$, refined $0.815$).

  \item \textbf{Legacy contamination is exposed.} Restricting to the top 50\% by legacy
        magnitude drops legacy $\rho$ to $0.620$ while raising refined $\rho$ to $0.867$
        on the same pairs, confirming that contamination drives the largest legacy scores.

  \item \textbf{The centroid trajectory method achieves 100\% benchmark recall.}
        All 11 literature-curated cascade pairs are correctly recalled across all 10 seeds,
        compared to partial recall by both static methods. This is the strongest
        validation result to date.

  \item \textbf{Pre-registered controls behave as expected.} IL-12 $\to$ IFN-$\gamma$
        ranks consistently in the upper decile with correct directionality; IL-6/IL-10
        is correctly reclassified as a shared-pathway pair after PBS-RC.

  \item \textbf{IL-27 and IFN-$\beta$ emerge as robust cascade sources.} Their top-pair
        lists overlap substantially between independent batches and match documented
        downstream biology.
\end{enumerate}

\end{document}
